%!TEX root = ../main.tex
% APÉNDICE E — Defensa del ticket promedio (relevamiento de precios en Superprof)

\chapter{Defensa del precio: relevamiento de tarifas en Superprof}\label{apx:precio}

\begin{guia}
Respalda el ticket promedio del Cap.~\ref{cap:economica}. Las cifras del
relevamiento son las informadas por el autor; las conversiones a dólares
dependen de un tipo de cambio que debe fijarse antes de incorporar estos
valores al capítulo.
\end{guia}

\section{Relevamiento}

Se relevaron los avisos de apoyo escolar publicados en el sitio de
Superprof Argentina \citep{superprof}, registrando el precio por hora que
cada profesor declara. El resultado fue:

\begin{table}[H]
    \centering
    \caption{Resultado del relevamiento (pesos argentinos por hora)}
    \label{tab:relevamiento-superprof}
    \small
    \begin{tabularx}{\textwidth}{L r}
        \toprule
        \tb{Indicador} & \tb{Valor} \\
        \midrule
        Avisos relevados                    & 414 \\
        Suma de precios horarios            & \$5.544.467 \\
        Precio horario promedio             & \$13.392,43 \\
        Mediana                             & \$12.000 \\
        Primer y tercer cuartil             & \$8.000 y \$18.000 \\
        Mínimo y máximo                     & \$1.990 y \$60.000 \\
        Avisos en modalidad online          & 387 (93\,\%) \\
        \bottomrule
    \end{tabularx}
\end{table}

El detalle de los 414 avisos se conserva en
\texttt{apendices/superprof-precios.csv}. Este promedio funciona como precio de referencia de mercado: es lo que
hoy cobra un profesor particular en la plataforma líder, contra lo cual
compite la oferta de \sistema{}.

\section{Traslado al ticket promedio}

El ticket del modelo es un promedio ya ponderado por la mezcla de
duraciones (Ec.~\ref{eq:duracion}, $\bar{d}=51$ minutos), por lo que el
precio horario se multiplica por $51/60 = 0{,}85$ antes de usarse. Con
el margen unitario $m(T) = T\,(0{,}27 - 0{,}0604) - 0{,}0255$, la
tabla~\ref{tab:sens-precio} muestra el efecto según el tipo de cambio,
manteniendo el resto de los supuestos del escenario base.

\begin{table}[H]
    \centering
    \caption{Efecto del nuevo precio de referencia (tipos de cambio ilustrativos)}
    \label{tab:sens-precio}
    \small
    \begin{tabularx}{\textwidth}{r r r r r r r}
        \toprule
        \tb{ARS/USD} & \tb{USD/h} & \tb{Ticket} & \tb{Margen} & \tb{Sesiones eq.\ Año 1} & \tb{Sesiones eq.\ Año 5} & \tb{VAN base} \\
        \midrule
        1.400 & 9,57 & 8,13 & 1,68 & 2.511 & 3.118 & 7.125 \\
        1.500 & 8,93 & 7,59 & 1,57 & 2.693 & 3.345 & 3.858 \\
        1.600 & 8,37 & 7,11 & 1,47 & 2.876 & 3.572 &   999 \\
        \midrule
        \multicolumn{2}{l}{Modelo vigente} & 7,00 & 1,44 & 2.924 & 3.631 & 308 \\
        \bottomrule
    \end{tabularx}
\end{table}

\section{Limitaciones}

\begin{itemize}
    \item Son precios \emph{pedidos} por los profesores, no pagados por los
          alumnos; los avisos pueden incluir descuentos o primeras clases
          gratis.
    \item 36 avisos figuran a \$1.990 por hora, probablemente un precio
          promocional o de entrada. Excluyéndolos, el promedio sube a
          \$14.478 y la mediana a \$14.000. Se adopta el promedio completo
          por ser el más conservador.
    \item El tipo de cambio es un parámetro del modelo, no un dato: la
          inflación en pesos erosiona el precio en dólares entre
          relevamientos.
\end{itemize}
